% Póster digital del proyecto final. Compilar desde esta carpeta: tectonic poster.tex
% Las figuras salen de scripts/graficas.py (resultados/graficas/*.pdf).
\documentclass[20pt, a0paper, landscape, margin=0mm, innermargin=18mm, blockverticalspace=6mm, colspace=18mm]{tikzposter}

\usepackage[spanish, es-noshorthands]{babel}
\usepackage{fontspec}
\usepackage{booktabs}
\usepackage{graphicx}
\usetikzlibrary{positioning, arrows.meta, calc}
\renewcommand{\familydefault}{\sfdefault}
\tikzposterlatexaffectionproofoff

% Azul marino del encabezado y los títulos, verde UIS para los acentos; los colores de datos son los de scripts/graficas.py.
\definecolor{marino}{HTML}{1B365D}
\definecolor{uis}{HTML}{8DB92E}
\definecolor{celeste}{HTML}{E8F0FA}
\definecolor{tinta}{HTML}{1F2933}
\definecolor{azul}{HTML}{0072B2}
\definecolor{verde}{HTML}{009E73}
\definecolor{gris}{HTML}{999999}

\usetheme{Default}
\colorlet{backgroundcolor}{white}
\colorlet{framecolor}{white}
\colorlet{titlefgcolor}{marino}
\colorlet{titlebgcolor}{white}
\colorlet{blocktitlebgcolor}{white}
\colorlet{blocktitlefgcolor}{marino}
\colorlet{blockbodybgcolor}{white}
\colorlet{blockbodyfgcolor}{tinta}
\colorlet{innerblocktitlebgcolor}{celeste}
\colorlet{innerblockbodybgcolor}{celeste}
\colorlet{innerblocktitlefgcolor}{marino}
\colorlet{innerblockbodyfgcolor}{tinta}

% Bloques sin marco: título azul alineado a la izquierda con una línea verde debajo.
\defineblockstyle{Linea}{
  titlewidthscale=1, bodywidthscale=1, titleleft, titleoffsetx=0pt, titleoffsety=0pt, bodyoffsetx=0pt,
  bodyoffsety=0pt, bodyverticalshift=0pt, roundedcorners=0, linewidth=0pt, titleinnersep=6mm, bodyinnersep=6mm
}{
  \ifBlockHasTitle
    \draw[draw=uis, line width=3mm] (blocktitle.south west) -- (blocktitle.south east);
  \fi
}
\useblockstyle{Linea}
\defineinnerblockstyle{Panel}{
  titlewidthscale=1, bodywidthscale=1, titleleft, titleoffsetx=0pt, titleoffsety=0pt, bodyoffsetx=0pt,
  bodyoffsety=0pt, bodyverticalshift=0pt, roundedcorners=12, linewidth=0pt, titleinnersep=5mm, bodyinnersep=7mm
}{
  \ifInnerblockHasTitle
    \draw[fill=innerblockbodybgcolor, draw=none, rounded corners=12] (innerblockbody.south west) rectangle (innerblocktitle.north east);
  \else
    \draw[fill=innerblockbodybgcolor, draw=none, rounded corners=12] (innerblockbody.south west) rectangle (innerblockbody.north east);
  \fi
}
\useinnerblockstyle{Panel}

\newcommand{\figura}[2]{\IfFileExists{#1}{\includegraphics[width=#2]{#1}}{\fbox{\parbox{#2}{\centering\vspace{4cm}Falta #1\vspace{4cm}}}}}
\newcommand{\dato}[1]{{\color{marino}\bfseries #1}}

% Estilos de los diagramas.
\tikzset{
  caja/.style={draw=marino, line width=2pt, rounded corners=6pt, fill=white, align=center, inner sep=10pt},
  flecha/.style={-{Latex[length=6mm]}, line width=3pt, draw=marino},
}

\title{¿Cuándo compensa la GPU para consultas analíticas?}
\author{Juan David Lipez Guevara \quad Janer Alberto Vega Jácome \quad Juan David Mena Gamboa \quad Camilo Carvajal Castro}
\institute{Programación Paralela, grupo 1 \quad Escuela de Ingeniería de Sistemas e Informática, Universidad Industrial de Santander\\Profesores: Carlos Jaime Barrios Hernández y Sergio Augusto Gélvez Cortés}

\makeatletter
\settitle{
  \begin{minipage}[c]{0.15\linewidth}\centering\includegraphics[width=0.8\linewidth]{logos/uis.png}\end{minipage}%
  \begin{minipage}[c]{0.7\linewidth}\centering\color{titlefgcolor}
    {\bfseries\fontsize{80}{90}\selectfont\@title\par}\vspace{0.3em}
    {\Large\color{marino!80}Tres consultas de TPC-H escritas a mano en serial, OpenMP y CUDA\par}\vspace{0.6em}
    {\LARGE\color{tinta}\@author\par}\vspace{0.4em}
    {\large\color{tinta}\@institute\par}
  \end{minipage}%
  \begin{minipage}[c]{0.15\linewidth}\centering
    \IfFileExists{logos/sc3.png}{\includegraphics[width=0.5\linewidth]{logos/sc3.png}}{}
  \end{minipage}
}
\makeatother

\begin{document}
\maketitle[width=1150mm, roundedcorners=0, linewidth=0pt, titletotopverticalspace=6mm, titletoblockverticalspace=4mm]

\begin{columns}

%%%%%%%%%%%%%%%% Columna 1: contexto y método %%%%%%%%%%%%%%%%
\column{0.27}

\block{Contexto}{
Las GPU dominan el entrenamiento de IA. ¿Ganan también en el trabajo diario de una base de datos analítica: filtrar, agregar y cruzar tablas de cientos de millones de filas?

\vspace{0.6em}
\innerblock{}{
\dato{Objetivo:} con el mismo código y los mismos datos, medir cuándo compensa cada arquitectura (un núcleo, una CPU de muchos núcleos o una GPU) y qué la limita.
}
}

\block{Qué construimos}{
Tres consultas escritas a mano en C++ y CUDA sobre datos guardados por columnas, cada una en tres versiones de la misma lógica y con verificación automática del resultado:

\vspace{0.6em}
\begin{center}
\begin{tikzpicture}
\node[caja] (gen) {Datos TPC-H\\{\small generados con DuckDB}};
\node[caja, right=1.6cm of gen] (col) {Columnas binarias\\{\small enteros en memoria}};
\node[caja, below=1.6cm of col, xshift=-6.2cm, fill=gris!15] (ser) {Serial\\{\small Volcano}};
\node[caja, right=0.8cm of ser, fill=azul!12] (omp) {OpenMP\\{\small morsels}};
\node[caja, right=0.8cm of omp, fill=verde!15] (gpu) {CUDA\\{\small tiles}};
\node[caja, below=1.6cm of omp] (val) {Resultado idéntico al de DuckDB\\{\small 495 de 495 mediciones}};
\draw[flecha] (gen) -- (col);
\foreach \b in {ser, omp, gpu} { \draw[flecha] (col.south) -- (\b.north); \draw[flecha] (\b.south) -- (val.north); }
\end{tikzpicture}
\end{center}
}

\block{Consultas de TPC-H}{
TPC-H es el benchmark estándar de bases de datos analíticas. Elegimos tres de sus 22 consultas, una por cada tipo de cuello de botella:

\vspace{0.4em}
\begin{itemize}
\item \dato{Q6:} filtra toda la tabla y suma el 2\,\% de las filas. \textit{Mide el ancho de banda de memoria.}
\item \dato{Q1:} agrega 300 millones de filas en 4 grupos. \textit{Mide la contención entre hilos.}
\item \dato{Q3:} cruza clientes, pedidos y líneas con una tabla hash. \textit{Mide los accesos aleatorios.}
\end{itemize}

\vspace{0.4em}
{\small SF 10: 60 millones de filas. SF 50: 300 millones.}
}

\block{Plataformas}{
\begin{itemize}
\item \dato{CPU:} nodo \texttt{smexa} del SC3-UIS, 2 AMD EPYC 9534 (128 núcleos, 2 nodos NUMA).
\item \dato{GPU:} NVIDIA T4 y A100 (Google Colab) y GTX Titan X (nodo \texttt{felix} del SC3-UIS).
\end{itemize}
}

\block{Dimensionamiento}{
\begin{itemize}
\item \textbf{Alcance:} 3 consultas y no las 22, una por cuello de botella; sin parser SQL ni optimizador.
\item \textbf{Datos:} SF 50 en CPU y SF 10 en GPU (cabe en los 16 GB de la T4); por eso se comparan filas por segundo.
\item \textbf{Clúster:} cuota de 1 GB y jobs de 1 hora; los datos se cargaron en el disco en RAM de cada nodo.
\end{itemize}
}

%%%%%%%%%%%%%%%% Columna 2: de serial a paralelo %%%%%%%%%%%%%%%%
\column{0.37}

\block{De serial a paralelo}{

\innerblock{1. Serial: modelo Volcano (Graefe, 1994)}{
\begin{minipage}[c]{0.34\linewidth}
\centering
\begin{tikzpicture}[node distance=1.1cm]
\node[caja, minimum width=7cm] (agg) {Agregación};
\node[caja, minimum width=7cm, below=of agg] (fil) {Filtro};
\node[caja, minimum width=7cm, below=of fil] (scan) {Scan \texttt{lineitem}};
\draw[flecha] (scan) -- node[right, font=\small] {1 fila} (fil);
\draw[flecha] (fil) -- node[right, font=\small] {1 fila} (agg);
\end{tikzpicture}
\end{minipage}%
\hspace{0.03\linewidth}%
\begin{minipage}[c]{0.62\linewidth}
Cada operador pide la siguiente fila a su hijo con una llamada \texttt{next()} virtual. Es el modelo clásico de los motores de bases de datos y nuestra línea base: un hilo, una fila a la vez.

\vspace{0.4em}
\dato{Costo:} una llamada indirecta por fila y por operador.
\end{minipage}
}

\vspace{0.8em}
\innerblock{2. CPU paralela: morsels con OpenMP (Leis et al., SIGMOD 2014)}{
\centering
\begin{tikzpicture}
\foreach \i in {0,...,11} { \pgfmathsetmacro{\tono}{10+6*\i} \draw[draw=marino, line width=1.5pt, fill=azul!\tono] (\i*1.9, 0) rectangle ++(1.9, 1.4); }
\node[above=2mm, font=\small] at (11.4, 1.4) {tabla dividida en morsels de 65\,536 filas};
\foreach \h/\x in {1/2.85, 2/8.55, 3/14.25, 4/19.95} {
  \node[caja, minimum width=4.8cm] (h\h) at (\x, -3.4) {Hilo \h\\{\small estado propio}};
  \draw[flecha] (\x, -0.1) -- (h\h.north);
}
\node[caja, fill=azul!15, minimum width=22cm] (comb) at (11.4, -7) {Combinar una sola vez al final};
\foreach \h in {1,...,4} { \draw[flecha] (h\h.south) -- (comb.north -| h\h.south); }
\end{tikzpicture}

\raggedright
\vspace{0.5em}
\begin{itemize}
\item Los hilos toman morsels de forma dinámica (\texttt{schedule(dynamic)}): ninguno queda ocioso si otro va lento.
\item Q6 y Q1: sumas y grupos locales por hilo, sin ramas y vectorizados por el compilador; cero atómicos compartidos.
\item Q3: tabla hash de direccionamiento abierto, construida por todos los hilos a la vez con \texttt{compare\_exchange}, e inicializada en paralelo para repartirla entre los dos sockets.
\end{itemize}
}

\vspace{0.8em}
\innerblock{3. GPU: tiles con CUDA (Shanbhag et al., SIGMOD 2020)}{
\centering
\begin{tikzpicture}
\node[caja] (ram) at (2.5, 0) {Memoria de la CPU};
\node[caja, fill=verde!10] (vram) at (17.5, 0) {Memoria de la GPU};
\draw[flecha] (ram) -- node[above, font=\small] {copia por PCIe} (vram);
\foreach \b in {0,...,3} {
  \node[draw=verde, line width=2pt, rounded corners=4pt, minimum width=5cm, minimum height=3.2cm, fill=white] (b\b) at (2.3+\b*5.8, -4.2) {};
  \node[font=\small, anchor=north, align=center] at (b\b.north) {Bloque \b: 512 filas\\128 hilos};
  \foreach \t in {0,...,7} { \draw[draw=verde, fill=verde!35] ($(b\b.south west)+(0.35+\t*0.56, 0.4)$) rectangle ++(0.42, 1.1); }
}
\node[caja, fill=verde!15, minimum width=21cm] (sum) at (11, -8.4) {Reducción en memoria compartida y \textbf{un} atómico global por bloque};
\foreach \b in {0,...,3} { \draw[flecha] (b\b.south) -- (sum.north -| b\b.south); }
\draw[flecha] (vram.south) -- (b3.north);
\end{tikzpicture}

\raggedright
\vspace{0.5em}
\begin{itemize}
\item Los filtros se evalúan como máscaras en registros; solo se leen las columnas de las filas que siguen vivas.
\item Q1: acumular en memoria compartida evita que miles de hilos escriban en las mismas 24 direcciones.
\item Q3: la tabla hash vive en la memoria de la GPU y se construye con \texttt{atomicCAS}.
\end{itemize}
}
}

\block{Lo que aprendimos}{
\begin{itemize}
\item La GPU compensa cuando los datos \textbf{ya viven en su memoria} y la consulta hace mucho trabajo por byte.
\item Antes de agregar hilos hay que saber \textbf{qué limita}: en Q6 era la memoria, no el cómputo.
\item Una inicialización serial y la memoria en un solo socket frenaban a Q3 (Amdahl y NUMA).
\end{itemize}
\textbf{Trabajo futuro:} copias desde memoria fijada, agregación por warp en Q1 y llevar la GPU a HIP para las AMD MI210 del SC3.
}

%%%%%%%%%%%%%%%% Columna 3: resultados %%%%%%%%%%%%%%%%
\column{0.36}

\block{Resultado: la GPU no gana siempre}{
\begin{center}
\figura{../resultados/graficas/central.pdf}{0.86\linewidth}
\end{center}
\begin{itemize}
\item \textbf{Con los datos ya en la GPU}, la A100 es 3,9 veces más rápida que 128 núcleos en el filtro y 8,5 veces en el join; la T4, una GPU de entrada, apenas empata.
\item \textbf{Copiando por PCIe}, todas las GPU pierden: la copia tarda hasta 100 veces más que el cálculo, también en la A100.
\item \textbf{En la agregación} gana la CPU, 2,7 veces frente a la A100 y 11 frente a la T4, por la contención de atómicos en la GPU.
\end{itemize}
}

\block{Escalado en CPU: el límite es la memoria}{
\begin{center}
\figura{../resultados/graficas/escalado.pdf}{\linewidth}
\end{center}
\begin{itemize}
\item \textbf{Q1 y Q3} escalan casi ideal hasta 32 hilos; Q3 llega a 52 veces con 128 hilos.
\item \textbf{Q6} sube por escalones: cada 8 núcleos comparten un enlace a memoria, y con 64 hilos ya lee al máximo de un socket.
\item \textbf{NUMA:} intercalar la memoria entre sockets hace a Q6 1,5 veces más rápida con 128 hilos.
\end{itemize}
}

\block{Contención en GPU y referencia}{
\begin{minipage}[c]{0.44\linewidth}
\figura{../resultados/graficas/contencion.pdf}{\linewidth}
\end{minipage}%
\hspace{0.03\linewidth}%
\begin{minipage}[c]{0.53\linewidth}
Acumular por bloque hace a Q1 \dato{9,4 veces} más rápida en la T4 y \dato{38 veces} en la A100.

\vspace{0.6em}
Frente a \dato{DuckDB}, en la misma máquina y con los mismos tipos, nuestras consultas escritas a mano son 1,5 veces más rápidas en el join y unas 5 veces en el filtro y la agregación: es la ganancia de especializar, no del paralelismo.
\end{minipage}
}

\end{columns}
\end{document}
